\section{Kernel of the adjoint Liouvillian}
\label{sec:liouvillian_kernel}
%% ---------------------------------------------------------

%% Source: \cite[Section~7.1, Theorem~7.2]{Wolf2012Quantum}.

\begin{definition}[Faithful stationary state]
    \label{def:has_faithful_stationary_state_ch12}
    \lean{HasFaithfulStationaryState}\leanok
    A generator $L$ has a faithful stationary state if there exists a positive definite
    density matrix $\rho_0$ with $L(\rho_0)=0$.
\end{definition}

\begin{definition}[Adjoint dissipator]
    \label{def:adjoint_dissipator_ch12}
    \lean{LindbladForm.adjointDissipator}\leanok
    \uses{def:lindblad_form}
    For a Lindblad operator $L_j$, the adjoint dissipator on observables is
    \[
        \mathcal{D}_{L_j}^*(A)
        = L_j^\dagger A L_j
        - \tfrac12 L_j^\dagger L_j A
        - \tfrac12 A L_j^\dagger L_j.
    \]
\end{definition}

\begin{definition}[Adjoint generator]
    \label{def:adjoint_generator_ch12}
    \lean{LindbladForm.toAdjointLinearMap}\leanok
    \uses{def:lindblad_form, def:adjoint_dissipator_ch12}
    The adjoint generator of a Lindblad form is
    \[
        L^*(A) = i[H,A] + \sum_j \mathcal{D}_{L_j}^*(A).
    \]
\end{definition}

\begin{theorem}[Trace-pairing characterization of the adjoint generator]
    \label{thm:trace_pairing_adjoint_generator_ch12}
    \lean{LindbladForm.trace_mul_toAdjointLinearMap_eq_trace_toLinearMap_mul}\leanok
    \uses{def:adjoint_generator_ch12, def:lindblad_form}
    For every density matrix $\rho$ and every observable $A$,
    \[
        \tr(\rho\,L^*(A)) = \tr(L(\rho)\,A).
    \]
\end{theorem}

\begin{proof}
    \leanok
    Expand the Schr\"odinger and Heisenberg formulas and apply cyclicity of the trace term by term.
\end{proof}

\begin{definition}[Commutant of the Lindblad data]
    \label{def:commutant_ch12}
    \lean{LindbladForm.commutant}\leanok
    \uses{def:lindblad_form}
    The commutant of a Lindblad form is the set of matrices commuting with $H$,
    with every Lindblad operator $L_j$, and with every adjoint $L_j^\dagger$.
\end{definition}

\begin{definition}[Adjoint kernel]
    \label{def:adjoint_kernel_ch12}
    \lean{LindbladForm.adjointKernel}\leanok
    \uses{def:adjoint_generator_ch12}
    The adjoint kernel is the set of observables $A$ satisfying $L^*(A)=0$.
\end{definition}

\begin{theorem}[Commutant lies in the adjoint kernel]
    \label{thm:commutant_subset_adjointKernel_ch12}
    \lean{LindbladForm.mem_adjointKernel_of_mem_commutant,
        LindbladForm.commutant_subset_adjointKernel}\leanok
    \uses{def:adjoint_generator_ch12, def:commutant_ch12, def:adjoint_kernel_ch12}
    If an observable commutes with $H$, with every $L_j$, and with every
    $L_j^\dagger$, then it lies in $\ker(L^*)$.
\end{theorem}

\begin{proof}
    \leanok
    The Hamiltonian commutator term vanishes, and each adjoint dissipator
    $\mathcal{D}_{L_j}^*(A)$ is zero under the commutation hypotheses.
\end{proof}

\begin{theorem}[Adjoint kernel lies in the commutant]
    \label{thm:adjointKernel_eq_commutant}
    \lean{LindbladForm.mem_commutant_of_mem_adjointKernel_of_hasFaithfulStationaryState}\leanok
    \uses{def:lindblad_form, def:has_faithful_stationary_state_ch12,
        def:commutant_ch12, def:adjoint_kernel_ch12}
    Let $L$ be a GKSL generator in Lindblad form with Hamiltonian $H$
    and Lindblad operators $\{L_j\}$.
    Define the adjoint generator $L^*$ by
    $\tr(\rho\,L^*(A)) = \tr(L(\rho)\,A)$.
    If $L$ has a faithful stationary state $\rho_0 > 0$ with $L(\rho_0) = 0$,
    then for any $A$:
    \[
        L^*(A) = 0
        \implies
        [A,H] = [A,L_j] = [A,L_j^\dagger] = 0\ \forall j.
    \]
    That is, $\ker(L^*) \subseteq \{H, L_j, L_j^\dagger\}'$.
    This is \cite[Theorem~7.2]{Wolf2012Quantum}.
\end{theorem}

\begin{proof}
    \leanok
    From $L^*(A) = 0$ and $L^*(A^\dagger) = 0$, compute
    $L^*(A^\dagger A) = \sum_j [A, L_j]^\dagger [A, L_j]$.
    Faithfulness of $\rho_0$ and $\tr(\rho_0 L^*(A^\dagger A)) = 0$
    force $[A, L_j] = 0$ for all $j$.
    Similarly $[A, L_j^\dagger] = 0$, and then $L^*(A) = 0$ directly gives
    $[A, H] = 0$.
\end{proof}

\begin{theorem}[Adjoint kernel equals the commutant]
    \label{thm:adjointKernel_eq_commutant_full}
    \lean{LindbladForm.adjointKernel_eq_commutant_of_hasFaithfulStationaryState}\leanok
    \uses{thm:adjointKernel_eq_commutant, def:commutant_ch12, def:adjoint_kernel_ch12}
    Under the same hypotheses as
    Theorem~\ref{thm:adjointKernel_eq_commutant},
    \[
        \ker(L^*) = \{H, L_j, L_j^\dagger\}',
    \]
    i.e.\ the adjoint kernel equals the commutant of $\{H, L_j, L_j^\dagger\}$.
    This is \cite[Theorem~7.2]{Wolf2012Quantum}.
\end{theorem}

\begin{proof}
    \leanok
    \uses{thm:adjointKernel_eq_commutant}
    The inclusion $\supseteq$ is the easy direction
    (commutant elements are in the kernel by a direct calculation).
    The reverse $\subseteq$ is
    Theorem~\ref{thm:adjointKernel_eq_commutant}.
\end{proof}

%% ---------------------------------------------------------
\section{Reducibility of quantum dynamical semigroups}
\label{sec:qds_reducibility}
%% ---------------------------------------------------------

%% Source: \cite[Section~7.1.2, Proposition~7.6]{Wolf2012Quantum}.

The four equivalent reducibility conditions of
\cite[Proposition~7.6]{Wolf2012Quantum} are as follows.

\begin{definition}[Nontrivial projection]
    \label{def:is_nontrivial_projection_ch12}
    \lean{IsNontrivialProjection}\leanok
    A projection is nontrivial if it is orthogonal and neither $0$ nor $\Id$.
\end{definition}

\begin{definition}[Rank-deficient fixed density]\label{def:has_rank_deficient_fixed_density}
    \lean{HasRankDeficientFixedDensity}\leanok
    \uses{def:is_nontrivial_projection_ch12, def:exp_semigroup_ch12}
    Condition~(1): there exists a density matrix $\rho_0$ with nontrivial
    kernel such that $e^{tL}(\rho_0) = \rho_0$ for all $t \ge 0$.
\end{definition}

\begin{definition}[Rank-deficient kernel element]\label{def:has_rank_deficient_kernel_element}
    \lean{HasRankDeficientKernelElement}\leanok
    \uses{def:is_nontrivial_projection_ch12}
    Condition~(2): there exists a density matrix $\rho_0$ with nontrivial
    kernel satisfying $L(\rho_0) = 0$.
\end{definition}

\begin{definition}[Invariant compression]\label{def:has_invariant_compression}
    \lean{HasInvariantCompression}\leanok
    \uses{def:is_nontrivial_projection_ch12, def:exp_semigroup_ch12}
    Condition~(3): there exists a nontrivial orthogonal projector $P$ such
    that $e^{tL}(P \MN{D} P) \subseteq P \MN{D} P$ for all $t \ge 0$.
\end{definition}

\begin{definition}[Block-upper-triangular Lindblad form]
    \label{def:has_block_upper_triangular_lindblad}
    \lean{HasBlockUpperTriangularLindblad}\leanok
    \uses{def:is_nontrivial_projection_ch12, def:lindblad_form}
    Condition~(4): there exists a nontrivial projector $P$ and a Lindblad form
    $(H, \{L_j\})$ for $L$ such that $(1-P)L_j P = 0$ and
    $(1-P)\kappa P = 0$ for all $j$, where
    $\kappa = iH + \frac{1}{2}\sum_j L_j^\dagger L_j$.
\end{definition}

\begin{definition}[Generator-preserved compression]
    \label{def:generator_preserves_compression_ch12}
    \lean{GeneratorPreservesCompression}\leanok
    \uses{def:has_invariant_compression}
    For a projection $P$, the generator preserves the compression $P M_d(\C) P$ if
    \[
        P\,L(PXP)\,P = L(PXP)
    \]
    for every $X \in M_d(\C)$.
\end{definition}

\begin{definition}[Reducible quantum dynamical semigroup]
    \label{def:is_reducible_qds_ch12}
    \lean{IsReducibleQDS}\leanok
    \uses{def:has_invariant_compression}
    A quantum dynamical semigroup is reducible if it has a nontrivial invariant compression.
\end{definition}

\begin{theorem}[Semigroup invariance implies generator invariance]
    \label{thm:generator_preserves_compression_from_semigroup_ch12}
    \lean{generatorPreservesCompression_of_semigroupPreservesCompression}\leanok
    \uses{def:generator_preserves_compression_ch12, def:has_invariant_compression}
    If the semigroup preserves $P M_d(\C) P$ for all non-negative times, then the generator
    preserves the same compression.
\end{theorem}

\begin{proof}
    \leanok
    Differentiate the identity $P e^{tL}(PXP) P = e^{tL}(PXP)$ at $t=0$.
\end{proof}

\begin{theorem}[Generator invariance implies semigroup invariance]
    \label{thm:semigroup_preserves_compression_from_generator_ch12}
    \lean{semigroup_preserves_compression_of_generator}\leanok
    \uses{def:generator_preserves_compression_ch12, def:has_invariant_compression}
    If the generator preserves $P M_d(\C) P$, then so does the exponential semigroup
    $e^{tL}$ for every $t \ge 0$.
\end{theorem}

\begin{proof}
    \leanok
    Apply the compression map termwise to the exponential series. Every iterate of $L$
    preserves the compression, so the whole series does as well.
\end{proof}

\begin{theorem}[Reducibility as a generator-level criterion]
    \label{thm:isReducibleQDS_iff_generator_preserves_compression_ch12}
    \lean{isReducibleQDS_iff_generator_preserves_compression}\leanok
    \uses{def:is_reducible_qds_ch12, def:generator_preserves_compression_ch12}
    A quantum dynamical semigroup is reducible if and only if its generator preserves some
    nontrivial compression.
\end{theorem}

\begin{proof}
    \leanok
    \uses{thm:generator_preserves_compression_from_semigroup_ch12,
        thm:semigroup_preserves_compression_from_generator_ch12}
    Combine Theorems~\ref{thm:generator_preserves_compression_from_semigroup_ch12}
    and~\ref{thm:semigroup_preserves_compression_from_generator_ch12}.
\end{proof}

\begin{theorem}[Rank-deficient fixed densities and kernel elements]
    \label{thm:wolf_prop_7_6_one_iff_two}
    \lean{wolf_prop_7_6_one_iff_two}\leanok
    \uses{def:has_rank_deficient_fixed_density,
        def:has_rank_deficient_kernel_element}
    A rank-deficient density matrix is a fixed point of $e^{tL}$ for all
    $t \ge 0$ if and only if it lies in $\ker L$.
    This is \cite[Proposition~7.6, (1)$\Leftrightarrow$(2)]{Wolf2012Quantum}.
\end{theorem}

\begin{proof}\leanok
    The equivalence $L(\rho) = 0 \iff e^{tL}(\rho) = \rho$ for all $t \ge 0$
    follows from the generator--semigroup correspondence: differentiate
    at $t=0$ for the forward direction, and exponentiate for the reverse.
\end{proof}

\begin{theorem}[A rank-deficient fixed density yields an invariant compression]
    \label{thm:wolf_prop_7_6_one_implies_three}
    \lean{wolf_prop_7_6_one_implies_three}\leanok
    \uses{def:has_rank_deficient_fixed_density,
        def:has_invariant_compression, def:is_gksl_generator}
    If $L$ is a GKSL generator and there exists a rank-deficient fixed
    density matrix, then the semigroup preserves a nontrivial compression.
    This is \cite[Proposition~7.6, (1)$\Rightarrow$(3)]{Wolf2012Quantum}.
\end{theorem}

\begin{proof}\leanok
    Let $\rho_0$ be a rank-deficient fixed density.
    Let $P$ be the support projection of $\rho_0$.
    Since $\rho_0$ is rank-deficient, $P \neq \Id$;
    since $\rho_0 \neq 0$, $P \neq 0$.
    Channel positivity and trace preservation force $e^{tL}$ to
    preserve $P \MN{D} P$.
\end{proof}

\begin{theorem}[Invariant compression yields block-upper-triangular Lindblad data]
    \label{thm:hasBlockUpperTriangularLindblad_of_hasInvariantCompression_ch12}
    \lean{hasBlockUpperTriangularLindblad_of_hasInvariantCompression}\leanok
    \uses{def:has_invariant_compression, def:has_block_upper_triangular_lindblad,
        def:is_gksl_generator, def:generator_preserves_compression_ch12}
    If a GKSL generator preserves a nontrivial compression, then its Lindblad data can be chosen
    block-upper-triangular with respect to that compression.
\end{theorem}

\begin{proof}
    \leanok
    \uses{thm:generator_preserves_compression_from_semigroup_ch12}
    First pass from semigroup invariance to generator invariance by
    Theorem~\ref{thm:generator_preserves_compression_from_semigroup_ch12}. The generator-level
    vanishing condition forces $(1-P)L_jP = 0$ and $(1-P)\kappa P = 0$ through the sum-of-squares
    identity.
\end{proof}

\begin{theorem}[Block-upper-triangular Lindblad data yield invariant compression]
    \label{thm:hasInvariantCompression_of_hasBlockUpperTriangularLindblad_ch12}
    \lean{hasInvariantCompression_of_hasBlockUpperTriangularLindblad}\leanok
    \uses{def:has_invariant_compression, def:has_block_upper_triangular_lindblad,
        def:generator_preserves_compression_ch12}
    If the Lindblad data are block-upper-triangular with respect to a nontrivial projection,
    then the associated semigroup preserves that compression.
\end{theorem}

\begin{proof}
    \leanok
    \uses{thm:semigroup_preserves_compression_from_generator_ch12}
    The block-upper-triangular relations imply generator-level compression invariance; then apply
    Theorem~\ref{thm:semigroup_preserves_compression_from_generator_ch12}.
\end{proof}

\begin{theorem}[Invariant compression and triangular Lindblad data]
    \label{thm:wolf_prop_7_6_three_iff_four}
    \lean{wolf_prop_7_6_three_implies_four, wolf_prop_7_6_four_implies_three}\leanok
    \uses{def:has_invariant_compression,
        def:has_block_upper_triangular_lindblad, def:is_gksl_generator}
    For a GKSL generator $L$, the semigroup preserves a nontrivial
    compression if and only if the Lindblad data can be chosen
    block-upper-triangular.
    This is \cite[Proposition~7.6, (3)$\Leftrightarrow$(4)]{Wolf2012Quantum}.
\end{theorem}

\begin{proof}\leanok
    \uses{thm:hasBlockUpperTriangularLindblad_of_hasInvariantCompression_ch12,
        thm:hasInvariantCompression_of_hasBlockUpperTriangularLindblad_ch12}
    Combine Theorems~\ref{thm:hasInvariantCompression_of_hasBlockUpperTriangularLindblad_ch12}
    and~\ref{thm:hasBlockUpperTriangularLindblad_of_hasInvariantCompression_ch12}.
\end{proof}

\begin{theorem}[Triangular Lindblad data give deficient kernels]
    \label{thm:blockUpperTriangular_implies_rankDeficient}
    \lean{hasRankDeficientKernelElement_of_hasBlockUpperTriangularLindblad}\leanok
    \uses{def:has_block_upper_triangular_lindblad,
        def:has_rank_deficient_kernel_element, def:is_gksl_generator}
    If the Lindblad data of a GKSL generator are block-upper-triangular
    with respect to some nontrivial projector $P$,
    then there exists a density matrix $\rho_0$ with nontrivial kernel
    satisfying $L(\rho_0) = 0$.
    This is \cite[Proposition~7.6, (4)$\Rightarrow$(2)]{Wolf2012Quantum}.
\end{theorem}

\begin{proof}
    \leanok
    The block-upper-triangular structure implies $L$ preserves the
    compression $P M_d(\C) P$.
    Therefore the semigroup $e^{tL}$ also preserves $P M_d(\C) P$
    for all $t \ge 0$.
    By the Brouwer fixed-point theorem, $e^{(1/m)L}$ has a fixed
    density matrix $\rho_m$ supported in $P M_d(\C) P$ for each $m$.
    Passing to a subsequential limit gives a density matrix
    $\rho_0$ with $P \rho_0 P = \rho_0$ and $L(\rho_0) = 0$.
    Since $P \neq \Id$, the matrix $\rho_0$ has nontrivial kernel.
\end{proof}

\begin{theorem}[Equivalent formulations of reducibility]
    \label{thm:wolf_prop_7_6_full}
    \lean{wolf_prop_7_6_full_equivalence}\leanok
    \uses{def:is_gksl_generator, def:has_rank_deficient_fixed_density,
        def:has_rank_deficient_kernel_element,
        def:has_invariant_compression,
        def:has_block_upper_triangular_lindblad}
    For a GKSL generator $L : \MN{D} \to \MN{D}$, the four conditions
    \begin{enumerate}
        \item rank-deficient fixed density,
        \item rank-deficient kernel element,
        \item invariant compression, and
        \item block-upper-triangular Lindblad form
    \end{enumerate}
    are equivalent.
    This is \cite[Proposition~7.6]{Wolf2012Quantum}.
\end{theorem}

\begin{proof}\leanok
    \uses{thm:wolf_prop_7_6_one_iff_two,
        thm:wolf_prop_7_6_one_implies_three,
        thm:wolf_prop_7_6_three_iff_four,
        thm:blockUpperTriangular_implies_rankDeficient}
    The cycle (1)$\leftrightarrow$(2), (1)$\to$(3),
    (3)$\leftrightarrow$(4), and (4)$\to$(2) closes
    the equivalence chain.
\end{proof}

%% Source: \cite[Corollary~7.2]{Wolf2012Quantum}.

\subsection{Sufficient conditions for non-reducibility}
\label{sec:cor_7_2_sufficient_conditions}

\begin{definition}[Lindblad span]\label{def:lindblad_span}
    \lean{lindbladSpan}\leanok
    \uses{def:lindblad_form}
    The $\C$-linear span of the Lindblad operators $\{L_j\}$.
\end{definition}

\begin{definition}[Hermitian-closed Lindblad span]\label{def:is_lindblad_span_hermitian_closed}
    \lean{IsLindbladSpanHermitianClosed}\leanok
    \uses{def:lindblad_span}
    The span of the Lindblad operators is closed under Hermitian conjugation,
    i.e.\ it forms a $*$-subspace of $M_d(\C)$.
\end{definition}

\begin{definition}[Trivial Lindblad commutant]\label{def:has_lindblad_span_trivial_commutant}
    \lean{HasLindbladSpanTrivialCommutant}\leanok
    \uses{def:lindblad_form}
    The commutant of the Lindblad family $\{L_j\}'$ equals $\C \cdot \Id$,
    i.e.\ the Lindblad operators act irreducibly on $M_d(\C)$.
\end{definition}

\begin{definition}[Kossakowski rank]\label{def:kossakowski_rank}
    \lean{kossakowskiRank}\leanok
    \uses{def:lindblad_form}
    The minimal number of Lindblad operators across all GKSL representations
    of $L$, which equals the rank of the Kossakowski matrix in the
    orthonormal-basis representation.
\end{definition}

\begin{theorem}[$\kappa$ block vanishing from generator compression]%
    \label{thm:kappa_block_generator_compression}
    \lean{kappa_block_of_generatorPreservesCompression}\leanok
    \uses{def:lindblad_form, def:has_block_upper_triangular_lindblad}
    Let $(H, \{L_j\})$ be a Lindblad form with
    $\kappa = iH + \tfrac{1}{2}\sum_j L_j^\dagger L_j$, and let $P$ be an
    orthogonal projection such that the generator preserves the compression
    $P M_d(\C) P$.
    If $(1-P) L_j P = 0$ for every $j$, then $(1-P) \kappa P = 0$.
\end{theorem}
\begin{proof}\leanok
    \uses{def:lindblad_form}
    From $L = \varphi - \kappa(\cdot) - (\cdot)\kappa^\dagger$ and
    $(1-P)L(P) = 0$: the $(1-P)\varphi(P)$ term vanishes because each
    $(1-P)L_j P = 0$, and the $(1-P)(P\kappa^\dagger)$ term vanishes
    because $(1-P)P = 0$. What remains is $(1-P)\kappa P = 0$.
\end{proof}

\begin{theorem}[Full algebra generation forbids block-upper-triangular form]%
    \label{thm:full_algebra_gen_no_block_ut}
    \lean{full_algebra_generation_implies_no_blockUpperTriangular}\leanok
    \uses{def:lindblad_form, def:has_block_upper_triangular_lindblad,
        thm:kappa_block_generator_compression}
    Let $(H, \{L_j\})$ be a Lindblad form.
    If the algebra generated by $\{L_j\}$ and $\kappa$ is the entire
    matrix algebra $M_d(\C)$, then the Lindblad data do not admit a
    block-upper-triangular decomposition.
\end{theorem}
\begin{proof}\leanok
    \uses{thm:kappa_block_generator_compression}
    Assume for contradiction that a block-upper-triangular decomposition
    exists with nontrivial projection $P$.
    By Theorem~\ref{thm:kappa_block_generator_compression},
    $(1-P)\kappa P = 0$.
    Induction on elements of the subalgebra shows
    $(1-P)A P = 0$ for every generator. Since the generated algebra is
    all of $M_d(\C)$, every matrix satisfies $(1-P)AP = 0$, forcing
    $P = 0$ or $P = \Id$ --- a contradiction.
\end{proof}

\begin{theorem}[Hermitian span with trivial commutant forbids triangular form]%
    \label{thm:herm_span_trivial_commutant_no_block_ut}
    \lean{hermitian_span_trivial_commutant_implies_no_blockUpperTriangular}\leanok
    \uses{def:lindblad_form, def:is_lindblad_span_hermitian_closed,
        def:has_lindblad_span_trivial_commutant,
        def:has_block_upper_triangular_lindblad}
    If the Lindblad span is closed under Hermitian conjugation and its
    commutant is $\C \cdot \Id$, then the Lindblad data do not admit a
    block-upper-triangular decomposition.
\end{theorem}
\begin{proof}\leanok
    Assume for contradiction that a block-upper-triangular decomposition
    exists with nontrivial projection $P$. The block vanishing
    $(1-P)L_j P = 0$ and Hermitian closure give $P L_j (1-P) = 0$ as
    well, so $[P, L_j] = 0$ for all $j$. But then $P$ lies in the
    commutant of the Lindblad span, contradicting triviality.
\end{proof}

\begin{theorem}[Large Kossakowski rank forbids block-upper-triangular form]
    \label{thm:large_kossakowski_rank_no_block_ut_ch12}
    \lean{large_kossakowski_rank_implies_no_blockUpperTriangular}\leanok
    \uses{def:kossakowski_rank, def:has_block_upper_triangular_lindblad}
    If the Kossakowski rank satisfies
    \[
        \rank(C) + d \ge d^2 + 1,
    \]
    then no block-upper-triangular Lindblad form exists.
\end{theorem}

\begin{proof}
    \leanok
    A block-upper-triangular traceless Lindblad family lies in a proper subspace whose dimension is
    at most $d^2-d$. The formal rank minimization argument then contradicts the large-rank bound.
\end{proof}

\begin{theorem}[No block-upper-triangular Lindblad form implies non-reducibility]
    \label{thm:not_isReducibleQDS_of_no_blockUpperTriangular_ch12}
    \lean{not_isReducibleQDS_of_no_blockUpperTriangular_lindblad}\leanok
    \uses{def:is_reducible_qds_ch12, def:has_block_upper_triangular_lindblad,
        thm:hasBlockUpperTriangularLindblad_of_hasInvariantCompression_ch12}
    If a GKSL generator admits no block-upper-triangular Lindblad form, then the associated
    quantum dynamical semigroup is not reducible.
\end{theorem}

\begin{proof}
    \leanok
    A reducible semigroup would have an invariant compression, and
    Theorem~\ref{thm:hasBlockUpperTriangularLindblad_of_hasInvariantCompression_ch12}
    would then produce a block-upper-triangular Lindblad form.
\end{proof}

\begin{theorem}[Full algebra generation implies non-reducibility]
    \label{thm:not_isReducible_of_generates_full_algebra}
    \lean{not_isReducible_of_generates_full_algebra}\leanok
    \uses{def:is_gksl_generator, def:lindblad_form,
        thm:full_algebra_gen_no_block_ut,
        thm:not_isReducibleQDS_of_no_blockUpperTriangular_ch12}
    If the algebra generated by $\{L_j\}$ and $\kappa$ is the entire
    matrix algebra $M_d(\C)$, then the QDS is not reducible.
    This is \cite[Corollary~7.2(1)]{Wolf2012Quantum}.
\end{theorem}
\begin{proof}\leanok
    \uses{thm:full_algebra_gen_no_block_ut,
        thm:not_isReducibleQDS_of_no_blockUpperTriangular_ch12}
    By Theorem~\ref{thm:full_algebra_gen_no_block_ut}, no block-upper-triangular
    decomposition exists. Then apply
    Theorem~\ref{thm:not_isReducibleQDS_of_no_blockUpperTriangular_ch12}.
\end{proof}

\begin{theorem}[Hermitian span and trivial commutant imply non-reducibility]
    \label{thm:not_isReducible_of_hermitian_span_trivial_commutant}
    \lean{not_isReducible_of_hermitian_span_trivial_commutant}\leanok
    \uses{def:is_gksl_generator, def:is_lindblad_span_hermitian_closed,
        def:has_lindblad_span_trivial_commutant,
        thm:herm_span_trivial_commutant_no_block_ut,
        thm:not_isReducibleQDS_of_no_blockUpperTriangular_ch12}
    If the Lindblad span is Hermitian-closed and its commutant is trivial,
    then the QDS is not reducible.
    This is \cite[Corollary~7.2(2)]{Wolf2012Quantum}.
\end{theorem}
\begin{proof}\leanok
    \uses{thm:herm_span_trivial_commutant_no_block_ut,
        thm:not_isReducibleQDS_of_no_blockUpperTriangular_ch12}
    By Theorem~\ref{thm:herm_span_trivial_commutant_no_block_ut}, no
    block-upper-triangular decomposition exists. Then apply
    Theorem~\ref{thm:not_isReducibleQDS_of_no_blockUpperTriangular_ch12}.
\end{proof}

\begin{theorem}[Large Kossakowski rank implies non-reducibility]
    \label{thm:not_isReducible_of_kossakowski_rank_ge}
    \lean{not_isReducible_of_kossakowski_rank_ge}\leanok
    \uses{def:is_gksl_generator, def:kossakowski_rank,
        thm:large_kossakowski_rank_no_block_ut_ch12,
        thm:not_isReducibleQDS_of_no_blockUpperTriangular_ch12}
    If $\rank(C) > d^2 - d$, then the QDS is not reducible.
    This is \cite[Corollary~7.2(3)]{Wolf2012Quantum}.
\end{theorem}
\begin{proof}\leanok
    The rank hypothesis rules out block-upper-triangular Lindblad data by
    Theorem~\ref{thm:large_kossakowski_rank_no_block_ut_ch12}. Then apply
    Theorem~\ref{thm:not_isReducibleQDS_of_no_blockUpperTriangular_ch12}.
\end{proof}
